\newcommand{\pluseq}{\mathrel{+\mkern-5mu=}}

\section{Automatizare}
\textit{Dificultate: 5/5} \\
\textit{Propunător: Ilie-Daniel Apostol} \\
\textit{Developer: Ilie-Daniel Apostol, Matteo-Alexandru Verzotti} \\
\textit{First solve: qfasfqwdfadafqwafvqvqvf după 02:07:04}  \\

Probabilitatea ca, la un pas oarecare al algoritmului, să obținem exact o operație reușită dorită de noi, este egală cu
\[
    q = \frac{p}{n - 1}
,\]

unde $\frac{1}{n - 1}$ este probabilitatea de a genera exact numărul dorit din intervalul $[1, n - 1]$, iar $p = \frac{P}{100}$ este probabilitatea ca operația să nu eșueze.

\medskip

Să ne gândim întâi la cele mai ``rele'' cazuri:
\begin{itemize}
    \item Elementul $1$ este mutat pe poziția $n$.
    Permutarea devine
    \[
        (2, 3, \ldots, n, 1).
    \]
    Pentru a o repara, operațiile trebuie efectuate
    în ordinea
    \[
        n-1 \;\longrightarrow\; n-2
        \;\longrightarrow\; \cdots
        \;\longrightarrow\; 1,
    \]
    iar fiecare trebuie să reușească. Între acestea pot
    apărea și alte operații, deoarece știm sigur că
    ele nu vor afecta permutarea.
\item Elementul $n$ este mutat pe poziția $1$.
    Permutarea devine
    \[
	(n, 1, 2, \dots, n - 1).
    \]
    Pentru a o repara, operațiile trebuie efectuate în ordinea
    \[
        1 \longrightarrow 2 \longrightarrow\ \dots \longrightarrow\ n - 1,
    \]
    
    cu aceleași condiții ca mai sus.
\end{itemize}

Vom nota cu $\uparrow$ șirul $(1, 2, \dots, n - 1)$, iar cu $\downarrow$ șirul $(n - 1, n - 2, \dots, 1)$.

Deoarece trebuie ca algoritmul nostru să sorteze permutarea în orice caz, înseamnă că acesta trebuie să conțină ca subșiruri atât pe $\uparrow$, cât și pe $\downarrow$. Se poate observa că un singur număr dintre acestea contribuie la ambele subșiruri, deci ne trebuie un algoritm ce va avea \textbf{minim} $2n - 3$ operații.

În continuare, vom calcula probabilitatea ca algoritmul nostru să le conțină pe $\uparrow$ și pe $\downarrow$ ca subșiruri folosind programare dinamică. Notăm cu $p_{t,a,b}$
probabilitatea ca, după $t$ pași, secvența să fi contribuit
cu $a$ numere la $\downarrow$ și cu $b$ numere la $\uparrow$. Pentru recurență, avem mai multe cazuri:

\begin{itemize}
    \item $a = b = n - 1$ \\
	Atunci deja am construit un algoritm bun, deci orice operație adăugată îl menține bun. Avem că:
    \[
	p_{t + 1, n - 1, n - 1} \pluseq p_{t, n - 1, n - 1}
    \]
\item $a = n - 1, b \neq n - 1$\\
    Atunci avem probabilitate $q$ ca următoarea operație generată să contribuie la $\downarrow$, și $1 - q$ să nu:
    \begin{align*}
	p_{t + 1, n - 1, b + 1} & \pluseq q \cdot p_{t, n - 1, b} \\
	p_{t + 1, n - 1, b} & \pluseq (1 - q) \cdot p_{t, n - 1, b}
    \end{align*}
\item $b = n - 1, a \neq n - 1$ \\
    Echivalent cu cazul de mai sus.
\item $a + 1 = n - b - 1$. \\
    Cu alte cuvinte, următoarea operație ar contribui atât la $\uparrow$, cât și la $\downarrow$. Ca și exemplu, dacă $n = 6, \; \uparrow = (1, 2), \; \downarrow = (5, 4)$, atunci generarea operației $(3, \textbf{reușită})$ contribuie la ambele șiruri. În acest caz, avem, la fel ca mai sus, probabilitate $q$ să contribuie la ambele și $1 - q$ la niciunul:
    \begin{align*}
	p_{t + 1, a + 1, b + 1} & \pluseq q \cdot p_{t, a, b} \\
	p_{t + 1, a, b} & \pluseq (1 - q) \cdot p_{t, a, b}
    \end{align*}
\item În restul cazurilor, avem probabilitate $q$ ca următoarea operație să contribuie la $\uparrow$, $q$ ca următoarea operație să contribuie la $\downarrow$, și $1 - 2q$ ca următoarea operație să nu contribuie la niciunul dintre ele.
    \begin{align*}
	p_{t + 1, a + 1, b} & \pluseq q \cdot p_{t, a, b} \\
	p_{t + 1, a, b + 1} & \pluseq q \cdot p_{t, a, b} \\
	p_{t + 1, a, b} & \pluseq (1 - 2q) \cdot p_{t, a, b}
    \end{align*}
\end{itemize}

Cazul de bază este $p_{0, 0, 0} = 1$, iar răspunsul se află în $p_{T, n - 1, n - 1}$. Pentru obținerea punctajului maxim, trebuie reținute doar ultimele $2$ linii din matricea $p$, pentru a nu depăși limita de memorie.

\medskip

Complexitate finală: $\mathcal{O}(TN^{2})$.
